\documentclass[11pt,a4paper]{article}
\usepackage[T1]{fontenc}
\usepackage[utf8]{inputenc}
\usepackage{amsmath,amsthm}
\usepackage{newtxtext,newtxmath}
\usepackage[margin=25mm,headheight=15pt]{geometry}
\usepackage{microtype}
\usepackage{mathtools}
\usepackage{booktabs,array,longtable}
\usepackage{xcolor}
\usepackage{enumitem}
\usepackage{fancyhdr}
\usepackage{titlesec}
\usepackage{tcolorbox}
\usepackage{xurl}
\definecolor{Ink}{HTML}{172D44}
\definecolor{Accent}{HTML}{176B72}
\definecolor{Pale}{HTML}{F0F5F7}
\usepackage[colorlinks=true,linkcolor=Ink,citecolor=Accent,urlcolor=Accent,
  pdftitle={Prime-supported numerators of arithmetic Hankel determinants},
  pdfauthor={AI-assisted research note},pdfsubject={OEIS A122249 and A122251}]{hyperref}
\titleformat{\section}{\large\bfseries\color{Ink}}{\thesection}{0.65em}{}
\titleformat{\subsection}{\normalsize\bfseries\color{Accent}}{\thesubsection}{0.65em}{}
\pagestyle{fancy}
\fancyhf{}
\fancyhead[L]{\small\color{Ink}Arithmetic Hankel determinants}
\fancyhead[R]{\small\color{Accent}Research note}
\fancyfoot[C]{\small\thepage}
\renewcommand{\headrulewidth}{0.35pt}
\setlength{\parskip}{4pt}
\setlength{\parindent}{0pt}
\setlength{\emergencystretch}{2em}
\setlist{nosep,leftmargin=1.6em}
\newtheorem{theorem}{Theorem}[section]
\newtheorem{lemma}[theorem]{Lemma}
\newtheorem{proposition}[theorem]{Proposition}
\newtheorem{corollary}[theorem]{Corollary}
\theoremstyle{definition}
\newtheorem{definition}[theorem]{Definition}
\newtheorem{example}[theorem]{Example}
\theoremstyle{remark}
\newtheorem{remark}[theorem]{Remark}
\newcommand{\Z}{\mathbb Z}
\newcommand{\N}{\mathbb N}
\newcommand{\Q}{\mathbb Q}
\newcommand{\num}{\operatorname{num}}
\newcommand{\den}{\operatorname{den}}
\newcommand{\sfac}{F}
\newcommand{\HH}{\mathcal H}
\newcommand{\ord}{\operatorname{ord}}
\newcommand{\dd}{\mathrm d}
\newcommand{\floor}[1]{\left\lfloor #1\right\rfloor}
\newcommand{\abs}[1]{\left|#1\right|}
\newcommand{\freepart}[2]{\left[#1\right]_{\perp #2}}
\newcommand{\supported}[2]{\left[#1\right]_{#2}}
\begin{document}
\thispagestyle{empty}
{\small\bfseries\color{Accent}OEIS CONJECTURE ATTACK \hfill 18 SEPTEMBER 2026\par}
\vspace{8mm}
{\LARGE\bfseries\color{Ink}Prime-supported numerators of\\[2mm]
arithmetic Hankel determinants\par}
\vspace{4mm}
{\large A proof of the formulas recorded in OEIS A122249 and A122251,\\
with exact denominator valuations and a sharp gcd theorem\par}
\vspace{5mm}
{\small AI-assisted research note. Complete English proofs and exact computational artifacts.\par}
\vspace{6mm}
\begin{tcolorbox}[colback=Pale,colframe=Accent,boxrule=0.5pt,arc=1mm,
  left=3mm,right=3mm,top=2mm,bottom=2mm]
\textbf{Abstract.}
For positive integers $d,b$ with $\gcd(d,b)=1$, let
\[
 D_m(d,b)=\det_{0\le i,j<m}\frac1{b+d(i+j)},\qquad
 \sfac_m=\prod_{j=0}^{m-1}j!.
\]
We prove that the numerator of $D_m(d,b)$ in lowest terms is
\[
 \boxed{\num D_m(d,b)=
 \prod_{\ell\mid d}\ell^{\,v_\ell(d)m(m-1)+2v_\ell(\sfac_m)}.}
\]
Here $\ell$ ranges over primes and $v_\ell$ denotes prime valuation.
This numerator is a square and is independent of the coprime shift $b$. For prime $d=p$ and $m=n+1$, the formula is exactly the conjecture
recorded in OEIS A122251 and its binary special case A122249.
The proof reduces cancellation to a finite count of pairs in residue classes.
We also determine every denominator valuation, prove that the gcd of all
coprime-shift denominators is the $d$-coprime part of the inverse Hilbert
determinant, and construct two shifts realizing that gcd. Further results
cover inverse-matrix denominators, unequal arithmetic steps, and noncoprime
parameters. Exact computations check the proofs' formulas independently of
the numerator prediction.
\end{tcolorbox}
\vspace{2mm}
\textbf{Outcome.} The selected formula is proved, not refuted. The
counterexamples in Section~\ref{sec:boundaries} concern only tempting
extensions of its hypotheses, not the OEIS conjecture itself.

\textbf{Priority qualification.} The retrieved OEIS entries still label the
claims as conjectural. This note supplies a self-contained proof of those
claims and extensions. It does not certify that no equivalent result has
previously appeared elsewhere; the underlying Cauchy determinant identity
is classical.
\vfill
{\small\color{Accent}Companion files: \texttt{code/arithmetic\_hankel.py},
\texttt{code/verify.py}, exact data, and an OEIS submission draft.}
\newpage
\tableofcontents
\newpage

\section{The selected problem and its status}
\label{sec:problem}

For a rational sequence $(a_k)_{k\ge0}$, use the Hankel-transform convention
\[
 h_n=\det(a_{i+j})_{0\le i,j\le n}.
\]
Thus the term indexed by $n$ is an $(n+1)\times(n+1)$ determinant, not an
$n\times n$ determinant. This distinction matters for all exponents below.

OEIS A122251, contributed by Paul Barry on 27 August 2006, records the
following general conjecture for prime $p$~\cite{oeis251}:
\begin{equation}
 \num\det_{0\le i,j\le n}\frac1{p(i+j)+1}
  =(p^2)^{S_p(n)},\qquad
 S_p(n)=\sum_{k=1}^n\sum_{a=0}^{n}\floor{\frac{k}{p^a}}.
 \label{eq:oeis}
\end{equation}
The entry itself displays the case $p=3$. The binary case is recorded in
A122249, which retains a request for a proof~\cite{oeis249}.

The internal records retrieved for this investigation carried revision tags
\texttt{\#2 Mar 30 2012} for A122251 and
\texttt{\#6 Oct 27 2024} for A122249. These are the entry revision tags,
not the encyclopedia-wide update date printed in the page footer. Both
retrieved records retain the conjectural status language. This is evidence
for choosing the problem as a currently recorded OEIS conjecture; it is not
a comprehensive literature certification of global openness.

The obstruction in proving~\eqref{eq:oeis} is not evaluating the determinant
as a rational product. Cauchy's determinant formula does that immediately.
The real question is whether all prime factors other than $p$ cancel from
the unreduced numerator. Checking only its $p$-adic valuation would leave
that question unanswered. Our proof explicitly treats every other prime.

The successful approach introduces a general positive step $d$ and shift
$b$. This makes the cancellation mechanism more visible, rather than
making the problem harder. The original conjecture follows as a special
case of a complete reduced-numerator formula.

\section{The general theorem and the generating-function setting}
\label{sec:main}

For $m\ge0$ and positive integers $d,b$, define
\begin{equation}
 H_m(d,b)=\left(\frac1{b+d(i+j)}\right)_{0\le i,j<m},
 \qquad D_m(d,b)=\det H_m(d,b),
 \label{eq:matrix}
\end{equation}
with $D_0(d,b)=1$. Write $D_m(d,b)=N_m(d,b)/B_m(d,b)$ in lowest terms,
where numerator and denominator are positive. Positivity will also follow
from the product evaluation below.

For a positive integer $u$, write
\[
 \supported{u}{d}=\prod_{\ell\mid d}\ell^{v_\ell(u)},
 \qquad
 \freepart{u}{d}=\frac{u}{\supported{u}{d}}.
\]
The products run over prime divisors of $d$. Thus
$\supported{u}{d}$ is the part supported on those primes, whereas
$\freepart{u}{d}$ is the largest divisor of $u$ coprime to $d$.
For $d=1$, these quantities are $1$ and $u$, respectively. Finally, set
\[
 \sfac_m=\prod_{j=0}^{m-1}j!,\qquad
 P_m(d,b)=\prod_{i,j=0}^{m-1}\bigl(b+d(i+j)\bigr).
\]
Empty products are one.

\begin{theorem}[Complete numerator and denominator]
\label{thm:main}
If $d,b\ge1$, $\gcd(d,b)=1$, and $m\ge0$, then
\begin{align}
 N_m(d,b)
   &=d^{m(m-1)}\supported{\sfac_m}{d}^{\,2}
     =\prod_{\ell\mid d}
       \ell^{v_\ell(d)m(m-1)+2v_\ell(\sfac_m)},
       \label{eq:num-main}\\
 B_m(d,b)
   &=\frac{P_m(d,b)}{\freepart{\sfac_m}{d}^{\,2}}\ \in\Z_{>0}.
       \label{eq:den-main}
\end{align}
In particular, $N_m(d,b)$ is independent of $b$, is a perfect square, and
has no prime divisor outside those of $d$. Moreover, $B_m(d,b)$ is coprime
to $d$.
\end{theorem}

The integrality in~\eqref{eq:den-main} is a substantive assertion. It is
proved below rather than assumed when writing the quotient. Once that
integrality is known, the two quantities are automatically coprime, because
the numerator is supported on primes dividing $d$ while the denominator is
not.

\begin{corollary}[The recorded OEIS conjectures]
\label{cor:oeis}
Formula~\eqref{eq:oeis} holds for every prime $p$ and every $n\ge0$.
Consequently it proves the displayed formulas in both A122249 and A122251.
\end{corollary}
\begin{proof}
Set $d=p$, $b=1$, and $m=n+1$ in Theorem~\ref{thm:main}.
Counting prime-power divisors of a factorial gives
\[
 v_p(k!)=\sum_{a\ge1}\floor{\frac{k}{p^a}}.
\]
Therefore
\begin{align*}
 v_p(N_{n+1}(p,1))
 &=n(n+1)+2\sum_{k=1}^n v_p(k!)\\
 &=2\sum_{k=1}^n\left(k+\sum_{a\ge1}\floor{\frac{k}{p^a}}\right)
 =2S_p(n).
\end{align*}
The sum over $a$ can be truncated at $n$, because every later term vanishes.
No other prime survives by the theorem.
\end{proof}

\subsection{Why these are generating-function determinants}
The sequence behind~\eqref{eq:matrix} has ordinary generating function
\begin{equation}
 A_{d,b}(z)=\sum_{k\ge0}\frac{z^k}{b+dk}
 =\int_0^1\frac{t^{b-1}}{1-zt^d}\,\dd t
 =\frac1b\,{}_2F_1\left(1,\frac bd;1+\frac bd;z\right),
 \qquad |z|<1.
 \label{eq:gf-moment}
\end{equation}
The integral follows by expanding the geometric series; the hypergeometric
identity follows from
$(b/d)_k/(1+b/d)_k=b/(b+dk)$. Also,
\[
 \frac1{b+d(i+j)}=\int_0^1 t^{b-1}t^{di}t^{dj}\,\dd t.
\]
Thus $H_m(d,b)$ is a Gram matrix of linearly independent monomials for a
positive weight. This explains positivity and places the problem within
moment sequences and Hankel transforms, rather than only matrix arithmetic.

\section{Cauchy's identity and the cancellation lemma}
\label{sec:proof}

\subsection{The algebraic evaluation}
\begin{lemma}[Cauchy's double alternant]
\label{lem:cauchy}
For distinct $x_i$, distinct $y_j$, and nonzero $x_i+y_j$,
\begin{equation}
 \det_{0\le i,j<m}\frac1{x_i+y_j}
 =\frac{\displaystyle\prod_{i<j}(x_j-x_i)(y_j-y_i)}
        {\displaystyle\prod_{i,j}(x_i+y_j)}.
 \label{eq:cauchy}
\end{equation}
\end{lemma}
\begin{proof}
After multiplication by the denominator on the right, the determinant is a
polynomial alternating separately in the $x$'s and in the $y$'s. Hence it
is divisible by both Vandermonde products. Its total degree is at most
$m^2-m$, precisely the degree of their product, so the quotient is a
constant. Multiply the claimed identity by $x_{m-1}+y_{m-1}$ and substitute
$x_{m-1}=-y_{m-1}$. On the determinant side only the last-row, last-column
cofactor remains. On the product side the cross factors cancel in pairs,
leaving the same identity of size $m-1$. The constant is therefore unchanged
when $m$ decreases by one. It equals one when $m=1$.
\end{proof}
This is a classical identity; see, for example, equation (2.7) in
Krattenthaler's survey~\cite{krattenthaler}. The short proof is included to
make the argument self-contained.

Taking $x_i=di$ and $y_j=b+dj$, and using
$\prod_{i<j}(j-i)=\prod_{j=0}^{m-1}j!=\sfac_m$, gives
\begin{equation}
 D_m(d,b)=\frac{d^{m(m-1)}\sfac_m^2}{P_m(d,b)}.
 \label{eq:raw}
\end{equation}
In particular, this is an unreduced expression with positive numerator and
denominator. The remaining issue is its exact cancellation.

\subsection{A finite combinatorial identity}
For a modulus $h\ge1$, write $m=qh+r$ with $0\le r<h$. For
$0\le c<h$, define
\[
 E_{h,r}(c)=\#\{(u,v):0\le u,v<r,\ u+v\equiv c\pmod h\}.
\]

\begin{lemma}[Residue-count excess]
\label{lem:excess}
For the preceding parameters,
\begin{equation}
 \#\{(i,j):0\le i,j<m,\ i+j\equiv c\pmod h\}
 -2\sum_{j=0}^{m-1}\floor{\frac jh}
 =hq+E_{h,r}(c)\ge0.
 \label{eq:excess}
\end{equation}
\end{lemma}
\begin{proof}
The residue $u$ occurs among $0,\ldots,m-1$ exactly
$q+\mathbf1_{\{u<r\}}$ times. Consequently the pair count is
\begin{align*}
 \sum_{u\bmod h}
 (q+\mathbf1_{\{u<r\}})
 (q+\mathbf1_{\{c-u\bmod h<r\}})
 =hq^2+2qr+E_{h,r}(c).
\end{align*}
On the other hand, the values $0,\ldots,q-1$ each occur $h$ times in the
list of quotients $\floor{j/h}$, and the value $q$ occurs $r$ times. Thus
\[
 2\sum_{j=0}^{m-1}\floor{\frac jh}=hq(q-1)+2qr.
\]
Subtracting gives~\eqref{eq:excess}. Its right side is nonnegative because
it is a sum of nonnegative integers.
\end{proof}

The crucial feature is uniformity in the residue $c$. There is no need to
know where the denominator progression hits zero modulo $h$; every residue
has enough occurrences to cancel the corresponding factorial contribution.

\subsection{Proof of the main theorem}
\begin{proof}[Proof of Theorem~\ref{thm:main}]
The case $m=0$ is immediate. Assume $m\ge1$ and fix a prime $\ell$.
If $\ell\mid d$, then $\ell\nmid b$ by coprimality, and hence
$\ell\nmid P_m(d,b)$. Formula~\eqref{eq:raw} gives
\begin{equation}
 v_\ell(D_m(d,b))=v_\ell(d)m(m-1)+2v_\ell(\sfac_m).
 \label{eq:bad-prime}
\end{equation}
This is nonnegative.

Now suppose $\ell\nmid d$. For each $a\ge1$, put $h=\ell^a$ and let
$c_h$ be the least nonnegative residue of $-b d^{-1}$ modulo $h$.
Then
\[
 h\mid b+d(i+j)\quad\Longleftrightarrow\quad i+j\equiv c_h\pmod h.
\]
Counting prime powers in the denominator product and in the factorials,
and then applying Lemma~\ref{lem:excess}, yields
\begin{align}
 v_\ell(D_m(d,b))
 &=\sum_{a\ge1}\left(
 2\sum_{j=0}^{m-1}\floor{\frac j{\ell^a}}
 -\#\{(i,j):\ell^a\mid b+d(i+j)\}\right)\notag\\
 &=-\sum_{a\ge1}\left(
 \ell^a\floor{\frac m{\ell^a}}
 +E_{\ell^a,\,m\bmod\ell^a}(c_{\ell^a})\right)
 \le0.
 \label{eq:good-prime}
\end{align}
All sums are finite: sufficiently large prime powers divide neither the
positive entry denominators nor any factorial factor.

Thus no prime outside $d$ can occur in the reduced numerator, whereas every
prime dividing $d$ has exactly the exponent in~\eqref{eq:bad-prime}.
This proves~\eqref{eq:num-main}. Dividing~\eqref{eq:raw} by this numerator
proves~\eqref{eq:den-main}, including its integrality. Its factors involve
no prime dividing $d$. Finally, $m(m-1)$ is even, so every exponent in the
numerator is even, proving the square assertion.
\end{proof}

\section{Exact denominator valuations}
\label{sec:den}

The proof gives more information than numerator cancellation. Define
\[
 T_r(s)=\max\{0,\ r-|s-(r-1)|\}.
\]
For integer $s$, this is the number of pairs in $\{0,\ldots,r-1\}^2$
with ordinary sum $s$, including the convention $T_0(s)=0$.
Because $0\le u+v\le2r-2<2h$, we have
\begin{equation}
 E_{h,r}(c)=T_r(c)+T_r(c+h),\qquad 0\le c<h.
 \label{eq:triangle}
\end{equation}

\begin{theorem}[Explicit local denominator]
\label{thm:den}
Under the hypotheses of Theorem~\ref{thm:main}, $v_\ell(B_m(d,b))=0$
for $\ell\mid d$. If $\ell\nmid d$, then
\begin{equation}
 v_\ell(B_m(d,b))=
 \sum_{a\ge1}\left[
 \ell^a\floor{\frac m{\ell^a}}
 +T_{r_a}(c_a)+T_{r_a}(c_a+\ell^a)\right],
 \label{eq:den-local}
\end{equation}
where $r_a=m\bmod\ell^a$ and
$c_a=(-bd^{-1})\bmod\ell^a$ is in $[0,\ell^a)$.
For $m\ge1$, the sum can be truncated once
$\ell^a>b+2d(m-1)$.
\end{theorem}
\begin{proof}
For $\ell\nmid d$, the reduced numerator has zero $\ell$-valuation.
Hence the denominator valuation is the negative of~\eqref{eq:good-prime}.
Apply~\eqref{eq:triangle}. Beyond the stated bound no entry denominator
is divisible by $\ell^a$, and the same bound is at least $m$; both parts
of the difference in~\eqref{eq:good-prime} therefore vanish.
\end{proof}

\begin{example}
For $m=3,d=3,b=1$, the raw product is
\[
 D_3(3,1)
 =\frac{3^6(0!1!2!)^2}{1\cdot4^2\cdot7^3\cdot10^2\cdot13}
 =\frac{729}{1783600}.
\]
The factor $4$ from the squared superfactorial cancels completely, while
the power $3^6$ does not cancel at all. For $d=6$ at the same size, the only
prime factor of $2!$ divides $d$, so the factor $4$ does not cancel and the
numerator is $6^6\cdot4=186624$.
\end{example}

\section{A sharp universal divisor and the exact gcd of denominators}
\label{sec:gcd}

Let
\begin{equation}
 \HH_m=\prod_{j=0}^{m-1}(2j+1)\binom{2j}{j}^{\!2}.
 \label{eq:hilbert}
\end{equation}
This is the reciprocal of the ordinary Hilbert determinant
$D_m(1,1)$ and is the sequence A005249~\cite{oeis5249}.
For completeness, dividing~\eqref{eq:raw} at sizes $j+1$ and $j$
with $d=b=1$ gives
\[
 \frac{D_{j+1}(1,1)}{D_j(1,1)}
 =\frac1{(2j+1)\binom{2j}{j}^2},
\]
which proves~\eqref{eq:hilbert} by multiplication.

\begin{lemma}[The least tail count]
\label{lem:min-tail}
For $0\le r<h$,
\[
 E_{h,r}(c)\ge\max(0,2r-h)
 \quad\text{for every }c,
 \qquad E_{h,r}(h-1)=\max(0,2r-h).
\]
\end{lemma}
\begin{proof}
In the cyclic group $\Z/h\Z$, let $A=\{0,\ldots,r-1\}$. Then
$E_{h,r}(c)=|A\cap(c-A)|$. Both sets have size $r$, so their intersection
has size at least $2r-h$, as well as being nonnegative. When $c=h-1$,
the only possible ordinary sum is $h-1$, and the pairs realizing it number
$\max(0,2r-h)$.
\end{proof}

\begin{theorem}[Exact gcd over all primitive shifts]
\label{thm:gcd}
For each $d\ge1$ and $m\ge0$,
\begin{equation}
 \boxed{\gcd_{\substack{b\ge1\\\gcd(b,d)=1}} B_m(d,b)
       =\freepart{\HH_m}{d}.}
 \label{eq:gcd}
\end{equation}
In particular, $\freepart{\HH_m}{d}$ divides every such denominator.
\end{theorem}
\begin{proof}
The case $m=0$ is immediate. For a prime $\ell\nmid d$, combine
Theorem~\ref{thm:den} and Lemma~\ref{lem:min-tail}. The result is
\begin{equation}
 v_\ell(B_m(d,b))\ge
 \sum_{a\ge1}\left[
 \ell^a\floor{\frac m{\ell^a}}+
 \max(0,2(m\bmod\ell^a)-\ell^a)\right]
 =v_\ell(\HH_m).
 \label{eq:universal-bound}
\end{equation}
The equality follows by applying the exact formula to $d=b=1$, when
$c_a=\ell^a-1$ achieves the minimum at every level. Primes dividing $d$
never occur in $B_m(d,b)$. This proves the divisibility statement.

It remains to show sharpness, rather than merely a lower bound. Fix
$\ell\nmid d$ and choose $A\ge1$ such that $\ell^A>2m-1$.
By the Chinese remainder theorem, there is a positive $b$ satisfying
\[
 b\equiv d\pmod{\ell^A},\qquad b\equiv1\pmod d.
\]
The second congruence ensures $\gcd(b,d)=1$. For every $0\le i,j<m$,
\[
 b+d(i+j)\equiv d(1+i+j)\pmod{\ell^A}.
\]
Since $1\le1+i+j\le2m-1<\ell^A$ and $\ell\nmid d$, neither side is
zero modulo $\ell^A$. Their $\ell$-valuations are therefore equal:
\[
 v_\ell\bigl(b+d(i+j)\bigr)=v_\ell(1+i+j).
\]
Subtracting the same factorial valuation in~\eqref{eq:raw} shows
$v_\ell(B_m(d,b))=v_\ell(\HH_m)$. Each prime not dividing $d$ therefore
attains the stated minimum for some admissible shift. Taking the minimum
valuation prime by prime proves the exact gcd formula.
\end{proof}

\begin{corollary}[Two shifts suffice, constructively]
\label{cor:two}
The gcd in~\eqref{eq:gcd} is already the gcd of the denominators at $b=1$
and at one explicitly constructible positive coprime shift $b=b_*$.
\end{corollary}
\begin{proof}
Let $\mathcal P$ be the finite set of primes dividing $B_m(d,1)$.
For each $\ell\in\mathcal P$, choose $A_\ell$ with
$\ell^{A_\ell}>2m-1$, and set
$M=\prod_{\ell\in\mathcal P}\ell^{A_\ell}$.
All these primes are coprime to $d$. Choose $b_*$ by
\[
 b_*\equiv d\pmod M,\qquad b_*\equiv1\pmod d.
\]
The preceding sharpness argument now applies simultaneously to every prime
of $B_m(d,1)$. Thus the gcd with $B_m(d,b_*)$ has exactly the valuations
in~\eqref{eq:gcd}. No other prime can enter that two-integer gcd.
\end{proof}

This is not a claim that one shift always has the smallest denominator
itself. It is a statement about divisibility and gcd. For example,
\[
 B_2(2,1)=45,\qquad B_2(2,17)=128877,
 \qquad\gcd(45,128877)=3=\freepart{\HH_2}{2}.
\]
The code also implements the Chinese-remainder construction, producing
exact two-shift certificates rather than relying on a finite search prefix.

\section{A second proof through inverse-matrix integrality}
\label{sec:inverse}

There is a complementary structural explanation of why only primes dividing
$d$ can survive in the determinant numerator. In fact, the entire inverse
matrix has denominators supported on those primes.

Write $\Z[1/d]$ for the rational numbers whose reduced denominators have
no prime divisor outside those dividing $d$.

\begin{lemma}[Rational binomial integrality away from the denominator]
\label{lem:binom}
If $\alpha\in\Z[1/d]$ and $k\ge0$, then
$\binom{\alpha}{k}\in\Z[1/d]$.
\end{lemma}
\begin{proof}
Represent $\alpha=u/v$ with every prime of $v$ dividing $d$.
If any factor in $\alpha(\alpha-1)\cdots(\alpha-k+1)$ is zero, the
assertion is immediate. Otherwise fix a prime $\ell\nmid d$.
For each power $h=\ell^a$, the integers $u-vt$ for $0\le t<k$
contain at least $\floor{k/h}$ multiples of $h$, because $v$ is invertible
modulo $h$ and every block of $h$ consecutive $t$'s contains exactly one.
The numerator product has valuation at least
$\sum_{a\ge1}\floor{k/\ell^a}=v_\ell(k!)$.
Division by $v^k k!$ cannot introduce a denominator factor $\ell$.
This holds for every $\ell\nmid d$.
\end{proof}

\begin{theorem}[Inverse entries]
\label{thm:inverse}
For all positive integers $d,b$ and all $m\ge1$,
$H_m(d,b)^{-1}$ has entries in $\Z[1/d]$.
No coprimality assumption on $d,b$ is needed for this assertion.
\end{theorem}
\begin{proof}
Put $\alpha=b/d$ and $G=(1/(\alpha+i+j))_{0\le i,j<m}$, so
$H_m(d,b)=d^{-1}G$. A cofactor calculation using
Lemma~\ref{lem:cauchy} gives
\begin{equation}
 (G^{-1})_{ij}
 =\frac{(-1)^{i+j}
   \displaystyle\prod_{k=0}^{m-1}(\alpha+i+k)
   \displaystyle\prod_{k=0}^{m-1}(\alpha+j+k)}
  {(\alpha+i+j)\,i!j!(m-1-i)!(m-1-j)!}.
 \label{eq:inverse-product}
\end{equation}
Indeed, removing row $j$ and column $i$ from the Cauchy determinant removes
the corresponding two strings of denominator factors, with their common
factor restored once. The Vandermonde ratios give the four factorials,
and the cofactor sign is $(-1)^{i+j}$.

Rearranging the finite products in~\eqref{eq:inverse-product} gives
\begin{align}
 (G^{-1})_{ij}
 &=(-1)^{i+j}(\alpha+i+j)
 \binom{m+\alpha+i-1}{m-j-1}
 \binom{m+\alpha+j-1}{m-i-1}\notag\\[-1mm]
 &\hspace{28mm}\cdot
 \binom{\alpha+i+j-1}{i}
 \binom{\alpha+i+j-1}{j}.
 \label{eq:inverse-binomial}
\end{align}
Every lower binomial argument is nonnegative, and every upper argument lies
in $\Z[1/d]$. Lemma~\ref{lem:binom} proves that each factor belongs to
that ring. Hence $H_m(d,b)^{-1}=dG^{-1}$ does as well.
\end{proof}

Since $\det H_m(d,b)^{-1}=1/D_m(d,b)$, the theorem alone proves that the
reduced numerator of $D_m(d,b)$ is supported on primes dividing $d$.
When $\gcd(d,b)=1$, its exact valuations at those primes then follow
from~\eqref{eq:raw}. This supplies a second route to the numerator theorem.
The residue-count proof remains useful because it also yields the exact
denominators and their sharp gcd.

\section{An extension to unequal arithmetic steps}
\label{sec:two-steps}

The cancellation phenomenon is not limited to symmetric Hankel matrices.
For positive integers $a,b,c$, define
\[
 C_m(a,b,c)=\det_{0\le i,j<m}\frac1{c+ai+bj}.
\]

\begin{theorem}[Two-step prime support and denominator divisibility]
\label{thm:two-step}
The reduced numerator of $C_m(a,b,c)$ has no prime divisor outside those
of $ab$. Its reduced denominator is divisible by
$\freepart{\HH_m}{ab}$.
\end{theorem}
\begin{proof}
Cauchy's identity gives
\[
 C_m(a,b,c)
 =\frac{(ab)^{\binom m2}\sfac_m^2}
 {\prod_{i,j}(c+ai+bj)}.
\]
Fix $\ell\nmid ab$, put $h=\ell^s$, and write $m=qh+r$.
Because multiplication by $a$ and by $b$ permutes residues modulo $h$,
the residues of $ai$ occur $q$ times each plus one extra occurrence on a
set $A$ of size $r$. Similarly the residues of $bj$ have an extra occurrence
on a set $B$ of size $r$. Thus the count of pairs satisfying
$c+ai+bj\equiv0\pmod h$ is
\[
 hq^2+2qr+|A\cap(-c-B)|.
\]
Subtracting the factorial contribution gives
$hq+|A\cap(-c-B)|\ge hq+\max(0,2r-h)$.
Summing over $s\ge1$ proves
$v_\ell(C_m(a,b,c))\le-v_\ell(\HH_m)$, just as in
\eqref{eq:universal-bound}. This excludes $\ell$ from the numerator and
proves the claimed denominator divisibility at every prime outside $ab$.
\end{proof}

This theorem deliberately does not assert a shift-independent formula for
primes dividing only one of $a,b$. The uniform residue argument applies at
primes where both steps are units; the common-step coprime-shift case is
special because its other primes never divide any entry denominator.

\section{Boundary cases and exact counterexamples to false extensions}
\label{sec:boundaries}

\subsection{Replacing a prime by a composite in the OEIS exponent fails}
The theorem permits composite $d$, but the original expression involving
powers $p^a$ cannot be generalized by simply replacing every $p$ by $d$.
For $d=4$, $n=2$, the naive expression predicts
\[
 (4^2)^{\sum_{k=1}^2\sum_{a\ge0}\floor{k/4^a}}
 =16^3=4096.
\]
In fact,
\[
 D_3(4,1)=\frac{16384}{52360425}.
\]
The missing factor is $4$, coming from $(2!)^2$: its prime is still $2$,
although no positive $k\le2$ reaches the first composite modulus $4$.
Formula~\eqref{eq:num-main}, using prime valuations rather than powers of a
composite base, gives the correct answer.

\subsection{Noncoprime parameters require a scaling correction}
Let $g=\gcd(d,b)$, $d_0=d/g$, and $b_0=b/g$. Then
\[
 H_m(d,b)=g^{-1}H_m(d_0,b_0),\qquad
 D_m(d,b)=g^{-m}D_m(d_0,b_0).
\]
If the primitive determinant is $N_0/B_0$ in lowest terms, then
\begin{equation}
 \num D_m(d,b)=\frac{N_0}{\gcd(N_0,g^m)},\qquad
 \den D_m(d,b)=\frac{g^m B_0}{\gcd(N_0,g^m)}.
 \label{eq:nonprimitive}
\end{equation}
These fractions are reduced because $\gcd(N_0,B_0)=1$.
This gives a complete formula for all positive $d,b$, not only primitive
pairs. For example,
\[
 D_2(2,2)=\frac1{48},\qquad
 D_3(4,2)=\frac{32}{496125}.
\]
The first numerator is not the primitive-step prediction $4$; the second
is not a square. Thus neither shift independence nor the square assertion
should be stated without the coprimality hypothesis.

\subsection{Arbitrary sparse minors need not have the same prime support}
Consecutive index intervals are also important. Even with the ordinary
Hilbert kernel, taking both row and column index sets to be $\{0,2\}$ gives
\[
 \det\begin{pmatrix}1&1/3\\1/3&1/5\end{pmatrix}=\frac4{45}.
\]
The ambient step is one, yet a prime survives in the numerator. The
Vandermonde differences of the sparse sets introduce additional factors.
This does not contradict Theorem~\ref{thm:two-step}: this particular minor
itself has row and column steps $a=b=2$, which correctly allow the prime $2$.

These examples delimit the established theorem. None refutes the original
prime-step conjecture, which is true.

\section{Exponent generating functions and fast evaluation}
\label{sec:algorithms}

The exponent sequences in the binary and ternary cases are themselves
recorded in A122247 and A122250~\cite{oeis247,oeis250}. We include their
general-prime derivation because it connects the result directly to
generating functions and produces a fast exact implementation.

Write $S_p(n)$ as in~\eqref{eq:oeis}, allowing the inner sum to run over all
$a\ge0$. Its first difference is
\[
 S_p(n)-S_p(n-1)=\sum_{a\ge0}\floor{\frac n{p^a}}.
\]
For any integer $q\ge1$,
\[
 \sum_{n\ge0}\floor{\frac nq}z^n
 =\frac{z^q}{(1-z)(1-z^q)}.
\]
This follows by expressing $\floor{n/q}$ as the number of positive
multiples of $q$ not exceeding $n$ and interchanging the two sums.
Consequently
\begin{equation}
 \mathcal S_p(z):=\sum_{n\ge0}S_p(n)z^n
 =\frac1{(1-z)^2}\sum_{a\ge0}\frac{z^{p^a}}{1-z^{p^a}}.
 \label{eq:lambert}
\end{equation}
Every coefficient receives only finitely many contributions, so the
identity is valid formally, as well as analytically for $|z|<1$.
Separating the $a=0$ term yields the functional equation
\begin{equation}
 (1-z)^2\mathcal S_p(z)
 =\frac z{1-z}+(1-z^p)^2\mathcal S_p(z^p).
 \label{eq:mahler}
\end{equation}
These formulas concern the exponent sequence; they are not generating
functions for the much faster-growing numerator sequence itself.

For exact computation at a large individual index, no coefficient expansion
and no determinant are needed. If $h=p^a$ and $m=q_a h+r_a$, then
\begin{equation}
 v_p(\sfac_m)=\sum_{a\ge1}
 \left[\frac{p^a q_a(q_a-1)}2+q_a r_a\right].
 \label{eq:fast}
\end{equation}
Only powers $p^a<m$ contribute. This uses $O(\log_p m)$ arithmetic
iterations after the prime is known. For composite $d$, factor $d$ and
apply~\eqref{eq:fast} once for each prime divisor; the implementation never
factors a large determinant numerator or denominator.

The iteration bound is not a constant-time or logarithmic bit-complexity
claim for printing the answer. For fixed $d>1$, the numerator contains
$d^{m(m-1)}$ and has quadratically many bits in $m$. Its large-integer
construction and output must pay that cost. The factorization of a varying
input $d$ is likewise a separate computational cost.

For an independently computed sequence of complete determinants, the
Cauchy product gives the ratio
\begin{equation}
 \frac{D_{j+1}(d,b)}{D_j(d,b)}
 =\frac{d^{2j}(j!)^2}
 {(b+2dj)\displaystyle\prod_{k=j}^{2j-1}(b+dk)^2}.
 \label{eq:ratio}
\end{equation}
The test suite uses this formula without invoking the predicted numerator,
then asks whether exact fraction reduction produces that numerator.

\section{Exact computational audit and reproducibility}
\label{sec:verification}

All tests use Python integers and \texttt{fractions.Fraction}; there is no
floating-point determinant computation. The mathematical proofs establish
all-parameter statements. The finite tests serve to catch indexing errors,
incorrect exponents, inverse formulas, implementation bugs, and falsely
stated boundary cases.

There are three independent levels of calculation. Direct rational Gaussian
elimination starts from the matrix entries and uses no determinant identity.
The Cauchy product and size-ratio evaluator reduce exact rational products
without using the numerator theorem. The predicted numerator is computed
only from prime factors of $d$ and factorial valuations. Agreement among
these is checked explicitly.

The executed test ranges and counts are shown in
Table~\ref{tab:checks}. Matrix size $m=0$ is included where stated. The
product-ratio range overlaps the direct range, so these counts should not
be interpreted as disjoint samples.

\begin{table}[htbp]
\centering
\small
\begin{tabular}{@{}p{.42\linewidth}r p{.40\linewidth}@{}}
\toprule
Check & Count & Range or interpretation\\
\midrule
Direct primitive determinants & 3,231 & $1\le d,b\le24$, $(d,b)=1$, $0\le m\le8$\\
Primitive product-ratio cases & 35,244 & $1\le d,b\le40$, $(d,b)=1$, $0\le m\le35$\\
Direct nonprimitive determinants & 847 & $1\le d,b\le18$, $(d,b)>1$, $0\le m\le6$\\
Residue-count identities & 34,255 & $2\le h\le32$, $0\le m\le64$, every residue\\
Denominator prime valuations & 20,240 & Primitive $d,b\le13$, $m\le7$, primes $<80$\\
Constructed two-shift gcds & 180 & $1\le d\le20$, $0\le m\le8$\\
Full inverse matrices & 720 & $1\le d,b\le12$, $1\le m\le5$\\
Inverse entry checks & 7,920 & Exact matrix product and denominator support\\
Two-step support cases & 4,374 & $1\le a,b,c\le9$, $0\le m\le5$\\
Direct two-step determinants & 2,916 & The preceding range restricted to $m\le3$\\
Superfactorial valuations & 605 & $m\le120$, $p=2,3,5,7,11$\\
Published initial terms & 14 & Seven each for the binary and ternary sequences\\
False-extension regressions & 5 & Exact examples from Section~\ref{sec:boundaries}\\
\bottomrule
\end{tabular}
\caption{All listed exact checks passed. Counts are also stored in
\texttt{data/verification\_summary.json}.}
\label{tab:checks}
\end{table}

The complete inverse test multiplies the matrix by the claimed inverse and
checks every entry of the identity matrix. The gcd tests construct the
second shift using Corollary~\ref{cor:two}, rather than merely taking a gcd
of a convenient short list of shifts. The local-valuation tests compare
prime-power counting with direct valuations of reduced denominators.

A few small cases are useful for immediate inspection:
\begin{center}
\begin{tabular}{@{}rrr@{\qquad}l@{}}
\toprule
$m$ & $d$ & $b$ & $D_m(d,b)$ in lowest terms\\
\midrule
2&2&1&$4/45$\\
3&2&1&$256/496125$\\
2&3&1&$9/112$\\
3&3&1&$729/1783600$\\
3&4&1&$16384/52360425$\\
3&6&1&$186624/971568325$\\
2&2&2&$1/48$\\
3&4&2&$32/496125$\\
\bottomrule
\end{tabular}
\end{center}

To reproduce the computational audit, run
\begin{quote}\ttfamily
python3 code/verify.py
\end{quote}
from the package directory, or invoke the script by its absolute path.
Python 3.10 or later is sufficient and no third-party modules are required.
The audit executed for this note used Python 3.13.5. The script writes the
summary, selected exact determinants, numerator tables, and two-shift
certificates to \texttt{data/}. The source can be rebuilt with the included
\texttt{build.sh} on a TeX installation with the packages named in the
preamble.

\section{Conclusion and scope of the result}

The prime-step numerator conjecture recorded in A122251 is true for all
primes and all indices. Its binary special case in A122249 is covered by
the same proof. More generally, coprime changes of the shift leave the
reduced numerator unchanged, and its prime factors and exponents are
completely determined by the arithmetic step and a superfactorial.

The strongest additional conclusion is the exact gcd theorem
\eqref{eq:gcd}, together with the constructive two-shift realization.
It separates the denominator's unavoidable Hilbert contribution from the
additional factors created by a particular shift. The inverse-matrix and
two-step results show that the cancellation phenomenon has a broader
arithmetic explanation.

The report provides ordinary mathematical proofs, not a proof-assistant
formalization or a peer-reviewed priority determination. No external
submission or OEIS edit has been made. The included submission text is a
draft for review.

\appendix
\section{Compact proof suitable for an OEIS-linked note}
\label{sec:compact}

Let $p$ be prime and $m=n+1$. Cauchy's formula gives
\[
 D=\det_{0\le i,j<m}\frac1{1+p(i+j)}
 =\frac{p^{m(m-1)}\prod_{j=0}^{m-1}(j!)^2}
 {\prod_{i,j=0}^{m-1}(1+p(i+j))}.
\]
The denominator is coprime to $p$, so the exact $p$-valuation of $D$ is
$m(m-1)+2\sum_{j<m}v_p(j!)=2S_p(n)$.

To rule out every other numerator prime $\ell$, let $h=\ell^a$ and
$m=qh+r$, $0\le r<h$. Because $p$ is invertible modulo $h$, the number
of pairs with $h\mid1+p(i+j)$ is the number with $i+j$ in one specified
residue class. Every residue among $0,\ldots,m-1$ occurs $q$ times, with
one extra occurrence on a set of $r$ residues. Thus the pair count is
$hq^2+2qr+E$ for some $E\ge0$. The contribution of $h$ to the squared
factorial product is $hq(q-1)+2qr$. Their difference, denominator minus
numerator, is $hq+E\ge0$. Sum over $a\ge1$.

Hence $v_\ell(D)\le0$ for every $\ell\ne p$. Its reduced numerator is
therefore $p^{2S_p(n)}$, exactly the claimed formula. This proves the
binary and ternary entries and the stated general-prime conjecture.

\begin{thebibliography}{9}
\bibitem{oeis251}
The OEIS Foundation Inc., \emph{A122251: Numerators of Hankel transform of
$1/(3n+1)$ (conjecture)}, entry contributed by Paul Barry, 27 August 2006.
\url{https://oeis.org/A122251}; internal record
\url{https://oeis.org/A122251/internal}. Consulted 18 September 2026,
America/Los\_Angeles.

\bibitem{oeis249}
The OEIS Foundation Inc., \emph{A122249: Numerators of Hankel transform of
$1/(2n+1)$}, entry contributed by Paul Barry, 27 August 2006.
\url{https://oeis.org/A122249}; internal record
\url{https://oeis.org/A122249/internal}. Consulted 18 September 2026.

\bibitem{krattenthaler}
C. Krattenthaler, \emph{Advanced Determinant Calculus}, S\'eminaire
Lotharingien de Combinatoire \textbf{42} (1999), Article B42q, 67 pp.
Equation (2.7), Cauchy's double alternant.
\url{https://arxiv.org/abs/math/9902004}; accessible full text
\url{https://arxiv.org/html/math/9902004v3}.

\bibitem{oeis5249}
The OEIS Foundation Inc., \emph{A005249: Determinant of inverse Hilbert
matrix}. \url{https://oeis.org/A005249}. Consulted 18 September 2026.

\bibitem{oeis247}
The OEIS Foundation Inc., \emph{A122247: Partial sums of A005187}, entry
contributed by Paul Barry, 27 August 2006.
\url{https://oeis.org/A122247}. Consulted 18 September 2026.

\bibitem{oeis250}
The OEIS Foundation Inc., \emph{A122250: Partial sums of A004128}, entry
contributed by Paul Barry, 27 August 2006.
\url{https://oeis.org/A122250}. Consulted 18 September 2026.
\end{thebibliography}
\end{document}
